\section{Unity, coercion, grace, and the hard legacy}\label{sec:aug-controversies}

\subsection{The Donatist division was not a footnote}

When Augustine became a priest, African Christianity was divided between
rival episcopal communions. Augustine and Possidius place compromised clergy,
episcopal legitimacy, sacramental integrity, churches, violence, and the
identity of the true Church inside the dispute, but they represent the
Catholic side. Full collation of the surviving conference record and opponent
dossiers is required before reconstructing Donatist self-understanding in
equal detail. Terms such
as \emph{schismatic} and \emph{circumcellion} are therefore reported as
polemical classifications rather than allowed to substitute for that archive.

Augustine first argued for persuasion, debate, and patient demonstration. His
anti-Donatist writings appealed to the worldwide communion, the distinction
between sacramental validity and fruitful unity, and the mixed character of
the Church before judgment. He participated in encounters with rival bishops,
including Proculeianus of Hippo. Violence and intimidation did occur, but
attribution in polemical sources requires case-by-case care. The beating of
Possidius and attacks on Catholic clergy were real to Augustine's circle;
Donatists likewise had reason to fear state force.

Imperial policy shifted toward suppression. The constitution now numbered
\emph{Codex Theodosianus} 16.5.38, issued on 12 February 405, expressly joined
Donatists with Manichaeans under an enforced program of Catholic unity. That
legal text does not explain every local action, but it establishes that
Augustine's pastoral argument developed inside a coercive imperial regime.
Augustine's own position changed.
By Letter 93 to Vincentius, usually dated 407/408, he defended legal pressure
because, he said, experience showed that coercion had brought some people to
gratitude and unity. In that argument he invoked Luke 14:23 through the Latin
imperative \emph{compelle intrare}. Letter 185 to Boniface, around 417, gives the mature defense: discipline
should aim at correction rather than vengeance, and the Church should prefer
salvation to a freedom that leaves persons in destructive error. Augustine
opposed execution and sometimes interceded for offenders. Those limits matter,
but they do not make coercion voluntary. Property penalties, exile, and state
enforcement burdened real communities.

An imperially convened conference at Carthage in June 411 brought Catholic and
Donatist bishops before the commissioner Marcellinus. The unique surviving
manuscript of the \emph{Gesta}, BnF lat.~1546, opens by naming Augustine among
seven Catholic representatives and seven Donatist representatives; its
incomplete three-session record survives through a ninth-century Lorsch copy.
Possidius treats the event as decisive for the Catholic campaign
(\emph{Life} 13). The manuscript and CSEL critical edition now control the
document's architecture, but this study does not turn an adversarial imperial
trial into neutral dialogue or a complete archive of Donatist experience.

\begin{studyframe}
\textbf{The judgment required.} It is anachronistic to call Augustine the
founder of the medieval Inquisition. It is equally misleading to make his
anti-execution pleas erase his theological defense of coercive law. He
developed that defense under pastoral, social, and imperial pressures, intended
correction rather than annihilation, and helped give later Christians a
scriptural rationale capable of uses far beyond his immediate limits.
\end{studyframe}

\subsection{Pelagius, inherited sin, and the priority of grace}

After the sack of Rome, Augustine entered a controversy that his own works
associate with Pelagius, Caelestius, and, later, Julian of Eclanum. This
revision has not collated their surviving fragments, the conciliar and Roman
dossiers, or a checked modern reconstruction of the controversy's external
phases. It therefore does not make “Pelagianism” a transparent summary of
every opponent's exact position or assign secure dates to each turn.

The biographically controlled center is Augustine's argument. In
\emph{Confessions} 10.29.40 he asks God to grant what God commands and then to
command what God wills. His anti-Pelagian works develop that priority of grace
through arguments about human solidarity in Adam, damaged desire, baptism,
freedom, perseverance, and predestination. The person genuinely wills and
acts, in Augustine's account, but every saving movement, including healed
freedom, begins in gift. The source audit names the works used; a specialist
edition-level review remains necessary before presenting a finer doctrinal
sequence or an opponent's position in the opponent's own voice.

In works against Julian, Augustine joins inherited sin and concupiscence to a
defense of marriage, procreation, and the created body. This biography records
that tension without reconstructing the uncollated exchange line by line.
Later Christian reception is not identical with every line of Augustine's
polemic, and this source set is not sufficient to decide how particular Greek
or Latin predecessors, African baptismal practice, or the Latin text of Romans
5:12 shaped each formulation. Those are review questions, not claims carried
by unseen modern books.

\subsection{Jews, pagans, and the Christian Roman order}

Augustine wrote in a society where imperial Christianity had acquired legal
power while traditional cults, Jewish communities, and rival Christian groups
remained present. His \emph{City of God} rejected the claim that abandonment
of the old gods caused Rome's sack and denied that any earthly commonwealth
could be the final city of God. This desacralized Roman success without
creating modern religious neutrality. Augustine accepted Christian imperial
government, supported restrictions in some religious conflicts, and imagined
law as capable of corrective service.

His treatment of Jews combined harmful theological subordination with a
notable argument against their annihilation. He read Jewish dispersion as a
providential witness to the Scriptures Christians received. This “witness
doctrine” could restrain forced disappearance while fixing living Jews inside
a Christian account of their unbelief. It should be described as an ambiguous
late-antique legacy, not praised as modern tolerance or reduced to the worst
later policies carried out in different circumstances.

Augustine addressed a society structured by wealth, poverty, and slavery. In
\emph{City of God} 19.15 he traced domination to sin while accepting servitude
within household order in the fallen world. The sources collated for this
revision do not support a fuller account of his interventions in kidnapping,
slave trading, or captive redemption, so no additional claim is made here. The
gap between theological diagnosis and social transformation remains part of
the hard legacy.

\subsection{Disagreement without a single final system}

Augustine wrote for more than forty years after baptism and revised his
language in response to new questions. A sentence from an early dialogue, a
Donatist letter, a late anti-Pelagian treatise, and a sermon to catechumens do
not automatically form one simultaneous system. In the \emph{Retractationes}
he revisited his books in sequence, noting errors, clarifications, and places
where a reader could misunderstand him. The title means a reconsideration or
review, not wholesale repudiation.

His self-review is a form of humility and also an act of authorial control.
He wanted readers to encounter a corrected corpus in a chosen order.
Possidius's \emph{Indiculus} catalogued books, letters, and sermons, helping
preserve a literary Augustine. Modern editions also identify bodies of material
commonly called the Divjak letters and Dolbeau sermons. Their edition-level
content was not collated for this biography; they are mentioned only to mark
that the surviving and edited corpus has not always been a closed list.

\evidenceaxis{Augustine came to defend legal coercion of Donatists as
corrective.}
{Letters 93.5.17--19 and 185.6.21--7.30; Possidius, \emph{Life} 9--14;
\emph{Codex Theodosianus} 16.5.38; and the 411 \emph{Gesta}. The full legal
sequence and complete proceedings still require specialist collation.}
{His opposition to capital punishment and his stated medicinal aim qualify,
but do not cancel, the coercive character. Donatist testimony and the unequal
imperial setting remain necessary counterweights.}
